\documentclass[../../main2.tex]{subfiles}

\begin{document}

	% ======================================================================
	% 骨架文件：小节标题与追问链为终版；正文由后续会话填充。
	% 本章是 Part IV 的前提章：先建分类判据（守恒性），再生成四通道。
	% 原书错误 1（四章并列）在此根治：后面四章每章开头标注「加入什么约束」。
	% ======================================================================

	\chapter{通道生成器：有压与无压}
	\label{ch:channels}

	\begin{motivation}
		上一章遗留的问题：贡献分解的 A、B 还是抽象字母——「因素」具体是什么？第5章埋下的分叉（有压/无压）还没有被当作分类判据正式使用。\\
		\textbf{核心动机}：建立全书第一个真正的分类原则——守恒性——并用它从「传播」出发生成四条通道。后面四章每一条，都是本章生成树上的一个节点。
	\end{motivation}

	\begin{logicmap}
	\begin{tikzpicture}[node distance=0.9cm, every node/.style={draw, rectangle, rounded corners, align=center}]
	\node (q) {传播的两种数学形态（§8.1）};
	\node (c) [below=of q] {守恒性 = 分类判据（§8.2）};
	\node (t) [below=of c] {四通道生成树（§8.3）};
	\node (m1) [below left=1.0cm and 2.2cm of t] {模因（无压）};
	\node (m2) [below=1.0cm of t] {经济（+稀缺）};
	\node (m3) [below right=1.0cm and 2.2cm of t] {政治（+身份）};
	\node (m4) [below=2.0cm of t] {法律（定责编码）};
	\draw[-{Stealth}] (q) -- (c);
	\draw[-{Stealth}] (c) -- (t);
	\draw[-{Stealth}] (t) -- (m1);
	\draw[-{Stealth}] (t) -- (m2);
	\draw[-{Stealth}] (t) -- (m3);
	\draw[-{Stealth}] (t) -- (m4);
	\end{tikzpicture}
	\end{logicmap}

	% ==================================================================
	\section{传播的两种数学形态}
	传播有两种数学形态，区别在于原体是否保留。

	第一种：复制（replication）。$A \to A+B$。你把一个想法告诉我，你手里还有那个想法，我手里也多了一个。原体不消耗，总量只增不减。笑话、知识、信仰——都走这条通道。一个段子被转发一万次，它还是一万个段子，不多不少。

	第二种：转移（transfer）。$A \to B$。你买一杯咖啡，钱离开你的口袋进入商家账户；你卖一股股票，那笔资金从你转移到买方。原体不再保留，总量锁死。这是经济通道的底层操作。

	为什么这个区分重要？因为两种形态服从不同的守恒律。复制没有守恒约束——信息可以无限增殖。转移有严格守恒——价值总量不变，只是换了持有者。后面四条通道，全部建立在这个分叉上：模因走复制，经济走转移，政治是两者的混合，法律给两者提供定量框架。
	% ==================================================================
	\section{守恒性作为分类判据}
	守恒性是全书第一个真正的分类原则。

	判据很简单：看传播过程中原体是否保留。保留 → 无压（复制）→ 模因通道。不保留 → 有压（转移）→ 经济通道。再加一层约束：如果交换中加入了身份关系（人情、义务、权力），就从经济进入政治。如果博弈结果需要被编码成确定数字，法律登场。

	这个判据不是人为划定的，而是从数学结构里长出来的。复制和转移在群论层面就是不同的操作：复制是单射（injective，原体保留），转移是满射（surjective，原体消失）。它们不能通过连续变形互相转化——就像你不能把一条绳子不断开就变成两个环。所以基于它们构建的通道也是正交的：你不会看到「半复制半转移」的中间态，只有叠加。

	\begin{readingnote}
		下次看到一个社会现象，先问一句：这个东西在传播时，原体保留吗？保留就是模因层面的事（无压），不保留就是资源层面的事（有压）。这一个判断就能把你带进正确的分析框架。
	\end{readingnote}
	% ==================================================================
	\section{四通道的生成树}
	从「传播」出发，每一步加一个约束，生成一条通道：

	传播（无约束）→ 模因通道。最底层，一切社会行为的基座。
	+ 稀缺约束 → 经济通道。传播必须付费，守恒律生效。
	+ 身份权重 → 政治通道。交换携带不可货币化的偏好。
	+ 定责编码 → 法律通道。博弈结果必须给出确定数字。

	这四条通道不是并列的四个独立系统，而是一棵生成树上的四个节点。每个节点都是前一个节点加了约束的子集。这个结构决定了后面四章的写法：每章开头标注「加入什么约束」，读者随时能回溯到树上自己的位置。
	% ==================================================================
	\section{本章小结与衔接}
	本章建立了全书的分类基础：传播的两种数学形态（复制 vs 转移），守恒性作为分类判据，以及从传播出发的四通道生成树。后面四章——模因、经济、政治、法律——是这棵树上的四个展开。每章开头会标注「加入什么约束」，保持与本章的对应关系。

	\paragraph*{逻辑衔接}
	第9章「通道一：模因（无压扩增）」：生成树的根节点，不加任何约束的纯传播。
\end{document}
